\section{Conclusion}
\label{sec:conclusion}

We have characterized the magnetic field decay of a yokeless hexapole
electromagnetic actuator using frequency-domain transfer function measurements
at 22 distances and 15 frequencies (330 measurements, 100\% valid).

The inverse-square model $H(x) = a^2/(x+b)^2$ yields
$b = 1978.9 \pm \SI{1.2}{\um}$ with $R^2 > 0.991$, confirming that:

\begin{enumerate}
\item The monopole model accurately describes the far-field behavior of a
single pole over the range $x \in [10, 2000]\;\si{\um}$.

\item The offset $b$ is frequency-independent, validating it as a geometric
property of the magnetic circuit.

\item For the yokeless configuration, $b \gg \rtip = \SI{5}{\um}$, indicating
that the flux is not concentrated at the pole tip.
\end{enumerate}

The derived physical parameters
($\kpole = 1.57 \times 10^{-7}$ Wb/A,
$\Ra = 3.19 \times 10^8$ A/Wb) and the force analysis show that the yokeless
actuator produces sub-piconewton forces ($F < \SI{0.021}{\pico\newton}$),
approximately $240\times$ weaker than the yoke-based design of
Menq~\cite{Zhang2010} ($b = \SI{405}{\um}$).

These results establish $b$ as the appropriate figure of merit for comparing
magnetic actuator geometries and confirm that a ferromagnetic yoke is essential
for generating the nanonewton-scale forces required for magnetic microbead
manipulation.
